\renewcommand{\arraystretch}{1.05}
\setlength{\parskip}{2pt}
% Extend only the figure into the margins: 19 cm on a 21 cm portrait A4 page.
\providecommand{\UCchange}[1]{\textcolor{blue}{#1}}
\newcommand{\InlineUCModel}[1]{%
  \par\noindent
  \begin{minipage}{\linewidth}
    \centering
    \makebox[\linewidth][c]{%
      \includegraphics[width=19cm,height=0.60\textheight,keepaspectratio]{sections/images/UC-#1.png}%
    }
    \captionsetup{hypcap=false}
    \captionof{figure}{UC-#1: detailed use-case model.}
    \label{fig:uc-model-#1}
  \end{minipage}
  \par
}
% ==================================================
% 2. USE-CASE SPECIFICATIONS
% ==================================================
\section{Use-case Models and Specifications}
This section describes the use cases of the
Medication Safety Assistant research prototype.
Each use case represents a user goal and an
observable result.
The use cases are organized by common user
workflows and supporting roles. Their order
does not imply a required execution sequence.
HCP reviews are simulated.
\UCchange{UC-19 is documented as a future feature and is not part of the current implementation scope.}
% ==================================================
% USE-CASE OVERVIEW
% ==================================================
\subsection*{Use-Case Overview}
\begin{longtable}{|
    p{0.13\textwidth}|
    p{0.53\textwidth}|
    p{0.24\textwidth}|
}
\hline
\textbf{ID} &
\textbf{Use-case Name} &
\textbf{Actor(s)}
\\ \hline
\endfirsthead
\hline
\textbf{ID} &
\textbf{Use-case Name} &
\textbf{Actor(s)}
\\ \hline
\endhead
UC-01
& Sign In
& All registered roles
\\ \hline
UC-02
& Enter and Confirm Medication Information
& Patient
\\ \hline
UC-03
& Check Medication Safety
& Patient
\\ \hline
UC-04
& Request HCP Review
& Patient
\\ \hline
UC-05
& Review Medication Safety Case
& Healthcare Professional (HCP); Patient / General User
\\ \hline
UC-06
& Set Up Medication Reminders
& Patient
\\ \hline
UC-07
& Manage Medication Reminders
& Patient
\\ \hline
UC-08
& Record Dose Status
& Patient
\\ \hline
UC-09
& View Adherence History
& Patient
\\ \hline
UC-10
& View Medication and Safety History
& Patient
\\ \hline
UC-11
& Ask a Medication Question
& Patient
\\ \hline
UC-12
& Generate Patient Health Summary
& Patient
\\ \hline
UC-13
& Share Patient Health Summary with HCP
& Patient
\\ \hline
\UCchange{UC-14}
& \UCchange{Review Medication Reference Data}
& Dataset Reviewer
\\ \hline
UC-15
& Manage User Accounts and Roles
& Administrator
\\ \hline
\UCchange{UC-16}
& \UCchange{Report Technical Issue}
& \UCchange{All registered roles}
\\ \hline
\UCchange{UC-17}
& \UCchange{Handle Technical Support Request}
& Technical Support
\\ \hline
\UCchange{UC-18}
& \UCchange{Manage Patient Health Profile}
& Patient / General User
\\ \hline
\UCchange{UC-19}
& \UCchange{Manage Personal Medication Cabinet (Future Feature)}
& Patient / General User
\\ \hline
\UCchange{UC-20}
& \UCchange{Manage Personal Account}
& Registered User (Patient / HCP)
\\ \hline
\end{longtable}
% UC-01 to UC-05 - updated against Use-Case Specification v1.0 (1 Nov 2026)
% Required packages in the MAIN document preamble:
% \usepackage{xcolor}
% \usepackage{graphicx}
% \usepackage{float}
% \usepackage{longtable}
% Changes relative to the old specification are highlighted in blue.
% Diagram images are PNG. The supplied use-case specification text is preserved.
% Store UC-01.png ... UC-20.png in sections/images/ relative to main.tex.
% ==================== UC-01 ====================
% Optional overview diagram: enabled once the future PNG is present.
\IfFileExists{sections/images/UC_00_Overview_Attached.png}{%
  \clearpage
  \begin{center}
    \includegraphics[width=\textwidth,height=0.94\textheight,keepaspectratio]{sections/images/UC_00_Overview_Attached.png}
  \end{center}
}{}
\clearpage
\subsection{Use-case: Sign In}
\label{uc:uc-01}
\InlineUCModel{01}
\begin{longtable}{|p{0.20\textwidth}|p{0.70\textwidth}|}
\hline
\textbf{Use case Name} & UC-01: Sign In \\ \hline
\textbf{Brief description} & A registered user signs in and receives access to the features permitted by their role. \\ \hline
\textbf{Actors} & Patient / General User; Healthcare Professional (HCP); Dataset Reviewer; Administrator; Technical Support \\ \hline
\textbf{Basic Flow} &
1. The user enters an account name and password.\par
2. The system checks the credentials and account status.\par
3. The system verifies the assigned role and permissions.\par
4. The system creates an authenticated session and opens the permitted dashboard.\par
\\ \hline
\textbf{Alternative Flows} &
\UCchange{\textbf{A1. Incorrect credentials (step 2)}}\par
\UCchange{The system shows an error without creating a session. The user may retry or cancel.}\par
\par
\UCchange{\textbf{A2. Inactive or blocked account (step 2)}}\par
\UCchange{The system denies access. No session is created.}\par
\par
\UCchange{\textbf{A3. Role permissions are unavailable (step 3)}}\par
\UCchange{The system denies access rather than creating an unrestricted session.}\par
\\ \hline
\textbf{Pre-conditions} & The user has a registered account and can open the sign-in page. \\ \hline
\textbf{Post-conditions} &
\textbf{Success:} \UCchange{A session is created with verified role permissions.}\par
\textbf{Failure:} No authenticated session is created. \\ \hline
\end{longtable}
% ==================== UC-02 ====================
\clearpage
\subsection{Use-case: Enter and Confirm Medication Information}
\label{uc:uc-02}
\InlineUCModel{02}
\begin{longtable}{|p{0.20\textwidth}|p{0.70\textwidth}|}
\hline
\textbf{Use case Name} & UC-02: Enter and Confirm Medication Information \\ \hline
\textbf{Brief description} & The patient enters medication information through a supported method and confirms the recorded details before further use. \\ \hline
\textbf{Actors} & Patient / General User \\ \hline
\textbf{Basic Flow} &
\UCchange{1. The patient provides medication details by manual entry, OCR image, or chatbot text.}\par
\UCchange{2. The system processes the input and clarifies the intended task when needed.}\par
\UCchange{3. The system displays the recorded details and marks any missing or uncertain fields.}\par
\UCchange{4. The patient corrects errors and explicitly confirms that the details match the source.}\par
\UCchange{5. The system saves the confirmed entries and any remaining information gaps.}\par
\\ \hline
\textbf{Alternative Flows} &
\UCchange{\textbf{A1. Input is unclear or incomplete (steps 2--3)}}\par
\UCchange{The system asks for missing details or marks unresolved fields. The patient may correct them before confirming.}\par
\par
\UCchange{\textbf{A2. Input processing fails (step 2)}}\par
\UCchange{The system offers a retry or another supported input method. No unconfirmed entry is released.}\par
\par
\UCchange{\textbf{A3. Patient does not confirm (step 4)}}\par
\UCchange{The input remains unconfirmed and cannot be used for safety assessment or reminder activation.}\par
\\ \hline
\textbf{Pre-conditions} & The patient is signed in and can access medication entry. \\ \hline
\textbf{Post-conditions} &
\textbf{Success:} \UCchange{Confirmed input is saved for later identification and supported checks; confirmation does not prove medicine identity or safety.}\par
\textbf{Failure:} Unconfirmed input is not used for safety checks or reminder activation. \\ \hline
\end{longtable}
% ==================== UC-03 ====================
\clearpage
\subsection{Use-case: Check Medication Safety}
\label{uc:uc-03}
\InlineUCModel{03}
\begin{longtable}{|p{0.20\textwidth}|p{0.70\textwidth}|}
\hline
\textbf{Use case Name} & UC-03: Check Medication Safety \\ \hline
\textbf{Brief description} & \UCchange{The patient receives supported medication risk findings, warnings, data limits, and next steps without automatic treatment changes.} \\ \hline
\textbf{Actors} & Patient / General User \\ \hline
\textbf{Basic Flow} &
1. The patient requests a check using confirmed medication information.\par
\UCchange{2. The system identifies assessable medicines and their active ingredients from supported reference data.}\par
\UCchange{3. The system checks supported drug--drug interactions, duplicate ingredients, and applicable contraindication risks where data allow.}\par
\UCchange{4. The system records findings, risk levels, and the coverage or limits of each check.}\par
\UCchange{5. The system shows warnings and next steps, and offers UC-04 when severe or unresolved risks need HCP review.}\par
6. The system saves the assessment, its status, and its time.\par
\\ \hline
\textbf{Alternative Flows} &
\UCchange{\textbf{A1. Medicine or reference data is unavailable (steps 2--3)}}\par
\UCchange{The system marks affected checks as unassessed, performs independent supported checks where possible, and does not claim the medicine is safe.}\par
\par
\UCchange{\textbf{A2. Severe interaction is found (step 3)}}\par
\UCchange{The system records a severe warning and holds new reminders for affected medicines pending HCP review. Existing reminders and treatment stay unchanged.}\par
\par
\UCchange{\textbf{A3. Assessment fails (steps 2--4)}}\par
\UCchange{The system reports an incomplete or failed check and retains any recorded warnings and reminder restrictions.}\par
\\ \hline
\textbf{Pre-conditions} & The patient is signed in and has at least one confirmed medication record. \\ \hline
\textbf{Post-conditions} &
\textbf{Success:} \UCchange{The saved assessment shows findings and limits. Severe risks keep affected new reminders on hold until an eligible, recorded HCP confirmation allows them.}\par
\textbf{Failure:} \UCchange{No incomplete or failed check is presented as safe; treatment and existing reminders remain unchanged.} \\ \hline
\end{longtable}
% ==================== UC-04 ====================
\clearpage
\subsection{Use-case: \UCchange{Arrange HCP Review}}
\label{uc:uc-04}
\InlineUCModel{04}
\begin{longtable}{|p{0.20\textwidth}|p{0.70\textwidth}|}
\hline
\textbf{Use case Name} & \UCchange{UC-04: Arrange HCP Review} \\ \hline
\textbf{Brief description} & \UCchange{The patient arranges a simulated HCP review of a saved safety case, with case-specific sharing permission and authorized assignment.} \\ \hline
\textbf{Actors} & Patient / General User \\ \hline
\textbf{Basic Flow} &
\UCchange{1. The patient opens a saved safety case that requires or is selected for HCP review.}\par
\UCchange{2. The system shows the review case, available findings, missing details, and case-sharing scope.}\par
\UCchange{3. If valid case permission is missing, the patient reviews the scope and authorizes sharing.}\par
\UCchange{4. The system reuses an active request or creates a pending request and assigns an eligible, available, authorized HCP when possible.}\par
\UCchange{5. The system displays the review status to the patient.}\par
\\ \hline
\textbf{Alternative Flows} &
\UCchange{\textbf{A1. Sharing is not authorized (step 3)}}\par
\UCchange{The case waits for permission. No HCP receives access, and no unauthorized review is started.}\par
\par
\UCchange{\textbf{A2. No eligible HCP is available (step 4)}}\par
\UCchange{The request remains Pending Assignment; the system does not report it as reviewed.}\par
\par
\UCchange{\textbf{A3. An active request already exists (step 4)}}\par
\UCchange{The system shows the existing request instead of creating a duplicate.}\par
\par
\UCchange{\textbf{A4. Request cannot be saved (step 4)}}\par
\UCchange{The system reports the error; no new access or assignment is granted.}\par
\\ \hline
\textbf{Pre-conditions} & \UCchange{The patient is signed in and has a saved medication safety case. Review access requires valid case-sharing permission.} \\ \hline
\textbf{Post-conditions} &
\textbf{Success:} \UCchange{A case is waiting for authorization, pending assignment, or assigned to an eligible HCP; the patient can view its status.}\par
\textbf{Failure:} \UCchange{No unauthorized case access or duplicate review is created; existing safety warnings and restrictions remain in place.} \\ \hline
\end{longtable}
% ==================== UC-05 ====================
\clearpage
\subsection{Use-case: Review Medication Safety Case}
\label{uc:uc-05}
\InlineUCModel{05}
\begin{longtable}{|p{0.20\textwidth}|p{0.70\textwidth}|}
\hline
\textbf{Use case Name} & UC-05: Review Medication Safety Case \\ \hline
\textbf{Brief description} & \UCchange{An authorized HCP reviews a shared safety case, exchanges information with the patient, and records a simulated review outcome.} \\ \hline
\textbf{Actors} & Healthcare Professional (HCP); Patient / General User \\ \hline
\textbf{Basic Flow} &
\UCchange{1. The HCP selects an assigned case with valid patient sharing permission.}\par
\UCchange{2. The system shows the permitted case details, safety warnings, available evidence, and case conversation.}\par
\UCchange{3. The HCP reviews the case and records an outcome and comments, or requests more information.}\par
\UCchange{4. The system saves the outcome, responsible HCP, time, and related replies.}\par
\UCchange{5. The patient can view the outcome and replies; a valid recorded schedule confirmation triggers a recheck of affected pending new reminders.}\par
\\ \hline
\textbf{Alternative Flows} &
\UCchange{\textbf{A1. More information is needed (step 3)}}\par
\UCchange{The HCP asks in the same case conversation. The patient replies there, and the case returns to the HCP queue.}\par
\par
\UCchange{\textbf{A2. Patient sends a follow-up (after step 5)}}\par
\UCchange{The system saves the reply in the same thread and returns the case to the authorized HCP queue.}\par
\par
\UCchange{\textbf{A3. Case remains unresolved (step 3)}}\par
\UCchange{The HCP records the reason. Required safety warnings and new-reminder holds remain active.}\par
\par
\UCchange{\textbf{A4. Access is invalid or saving fails (steps 1 or 4)}}\par
\UCchange{The system blocks the action or reports failure; no unrecorded outcome is treated as approval.}\par
\\ \hline
\textbf{Pre-conditions} & \UCchange{The HCP is signed in, assigned to the case, and authorized by the patient to view the shared details.} \\ \hline
\textbf{Post-conditions} &
\textbf{Success:} \UCchange{A traceable simulated outcome or case-conversation update is saved and visible to permitted participants.}\par
\textbf{Failure:} \UCchange{No unauthorized information is disclosed, and no treatment or reminder restriction is changed by an unrecorded decision.} \\ \hline
\end{longtable}
% ---------- 04. Reminders and adherence ----------
% UC-06 - Set Up Medication Reminders | from UC-07
\clearpage
\subsection{Use-case: Set Up Medication Reminders}
\label{uc:uc-06}
\InlineUCModel{06}
\begin{longtable}{|p{0.20\textwidth}|p{0.71\textwidth}|}
\hline
\endfirsthead
\multicolumn{2}{l}{\textbf{UC-06: Specification (continued)}} \\
\hline
\endhead
\textbf{Use case Name} & UC-06: Set Up Medication Reminders \\ \hline
\textbf{Brief description} & The patient activates reminders from an existing, confirmed medication schedule when the applicable safety rules allow it. \\ \hline
\textbf{Actors} & Patient / General User \\ \hline
\textbf{Basic Flow} &
1. The patient selects a medication with existing prescription or recorded HCP instructions.\par
2. The system displays the recorded schedule, its source, and relevant safety status.\par
3. The patient confirms that the schedule matches the existing instructions.\par
4. The system checks schedule verification, recorded risk warnings, and the new-reminder safety gate.\par
5. The system activates eligible reminders and shows their status.
\\ \hline
\textbf{Alternative Flows} &
\textbf{A1. Schedule is missing or unverified (steps 2--3)}\par
1. The system does not activate a reminder without a verified schedule.\par
2. The patient may correct the information or end the use case.
\\ \hline
\textbf{Alternative Flows} &
\textbf{A2. Patient proposes new timing (step 3)}\par
1. The system records a proposed change as Pending HCP Review.\par
2. It does not activate that proposed schedule before the required simulated confirmation.
\\ \hline
\textbf{Alternative Flows} &
\textbf{A3. Severe interaction affects the medicine (step 4)}\par
1. The system keeps affected new reminders Pending HCP Review.\par
2. Other eligible reminders may activate; existing reminders are not stopped automatically.
\\ \hline
\textbf{Alternative Flows} &
\textbf{A4. Safety check is unavailable (step 4)}\par
1. The system reports Unassessed rather than safe clearance.\par
2. The reminder remains restricted when the applicable review policy requires confirmation.
\\ \hline
\textbf{Alternative Flows} &
\textbf{A5. Browser notifications are unavailable (step 5)}\par
1. The system reports the notification limitation and still shows eligible schedules in the application.
\\ \hline
\textbf{Pre-conditions} &
The patient is signed in and has a saved
medication record for which reminder setup
can be requested.
 \\ \hline
\textbf{Post-conditions} &
\textbf{Success:} A permitted reminder setting
is saved, or a timing proposal is recorded
as Pending HCP Review.\par
\textbf{Failure:} No unconfirmed or failed change
is applied. The previous reminder settings
remain effective.
\\ \hline
\end{longtable}
% UC-07 - Manage Medication Reminders | from UC-08
\clearpage
\subsection{Use-case: Manage Medication Reminders}
\label{uc:uc-07}
\InlineUCModel{07}
\begin{longtable}{|p{0.20\textwidth}|p{0.71\textwidth}|}
\hline
\endfirsthead
\multicolumn{2}{l}{\textbf{UC-07: Specification (continued)}} \\
\hline
\endhead
\textbf{Use case Name} & UC-07: Manage Medication Reminders \\ \hline
\textbf{Brief description} & The patient changes permitted reminder settings or disables notifications without changing prescribed treatment. \\ \hline
\textbf{Actors} & Patient / General User \\ \hline
\textbf{Basic Flow} &
1. The patient opens an existing reminder and chooses to manage it.\par
2. The system shows the current schedule, its source, and notification settings.\par
3. The patient requests a notification change or a schedule update.\par
4. The system checks whether the change is allowed by the verified instructions and safety rules.\par
5. The patient confirms the request.\par
6. The system saves an allowed setting or records a proposed timing change as Pending HCP Review.
\\ \hline
\textbf{Alternative Flows} &
\textbf{A1. Patient disables notifications (step 3)}\par
1. The system prepares to disable
only the selected notification.\par
2. Continue to step 5 for confirmation.
The medication instructions remain unchanged.
\\ \hline
\textbf{Alternative Flows} &
\textbf{**A2. Timing differs from verified**
**instructions (step 4)**}\par
1. The system marks the proposed
change as requiring HCP review.\par
2. Continue to step 5 for patient
confirmation. The existing schedule
remains unchanged.
\\ \hline
\textbf{Alternative Flows} &
\textbf{A3. Patient cancels or update fails (steps 5--6)}\par
1. The system keeps the previous reminder settings and reports an error when relevant.
\\ \hline
\textbf{Pre-conditions} & The patient is signed in and owns an existing reminder. \\ \hline
\textbf{Post-conditions} & \textbf{Success:} A permitted reminder setting is saved, or a timing proposal is Pending review.\par
\textbf{Failure:} The prior reminder settings remain effective unless the patient explicitly disabled a notification. \\ \hline
\end{longtable}
% UC-08 - Record Dose Status | from UC-09
\clearpage
\subsection{Use-case: Record Dose Status}
\label{uc:uc-08}
\InlineUCModel{08}
\begin{longtable}{|p{0.20\textwidth}|p{0.71\textwidth}|}
\hline
\endfirsthead
\multicolumn{2}{l}{\textbf{UC-08: Specification (continued)}} \\
\hline
\endhead
\textbf{Use case Name} & UC-08: Record Dose Status \\ \hline
\textbf{Brief description} & The patient records whether a scheduled medication dose was Taken, Missed, or Skipped for adherence tracking. \\ \hline
\textbf{Actors} & Patient / General User \\ \hline
\textbf{Basic Flow} &
1. The patient selects a scheduled dose from the reminder or adherence view.\par
2. The system shows the medicine, scheduled time, and current dose status.\par
3. The patient selects Taken, Missed, or Skipped.\par
4. The patient confirms the entry.\par
5. The system saves the dose status and recorded time.
\\ \hline
\textbf{Alternative Flows} &
\textbf{A1. Existing entry is corrected (step 3)}\par
1. The system shows the previous status for comparison.\par
2. The patient confirms the correction; continue to step 5.
\\ \hline
\textbf{Alternative Flows} &
\textbf{A2. Patient cancels (step 4)}\par
1. The system keeps the previous entry and ends the use case.
\\ \hline
\textbf{Alternative Flows} &
\textbf{A3. Dose event is unavailable or save fails (steps 2 or 5)}\par
1. The system explains that the entry cannot be saved.\par
2. No unsupported dose record is created.
\\ \hline
\textbf{Pre-conditions} & The patient is signed in and has an accessible scheduled dose or adherence entry. \\ \hline
\textbf{Post-conditions} & \textbf{Success:} The confirmed dose status is saved and available in adherence history.\par
\textbf{Failure:} No dose entry is changed; no treatment recommendation is produced. \\ \hline
\end{longtable}
% UC-09 - View Adherence History | from UC-10
\clearpage
\subsection{Use-case: View Adherence History}
\label{uc:uc-09}
\InlineUCModel{09}
\begin{longtable}{|p{0.20\textwidth}|p{0.71\textwidth}|}
\hline
\endfirsthead
\multicolumn{2}{l}{\textbf{UC-09: Specification (continued)}} \\
\hline
\endhead
\textbf{Use case Name} & UC-09: View Adherence History \\ \hline
\textbf{Brief description} & The patient reviews recorded medication use over time. \\ \hline
\textbf{Actors} & Patient / General User \\ \hline
\textbf{Basic Flow} &
1. The patient opens Adherence History.\par
2. The system shows available scheduled doses and recorded statuses.\par
3. The patient may select a medicine
or time period to filter the history.\par
4. The system displays the available
records matching the selected filter,
or all permitted records if no
filter is selected.
\\ \hline
\textbf{Alternative Flows} &
\textbf{A1. No entries match (step 4)}\par
1. The system shows an empty result and allows another filter.
\\ \hline
\textbf{Alternative Flows} &
\textbf{A2. History is incomplete or unavailable (step 2)}\par
1. The system displays available records and reports any unavailable period.
\\ \hline
\textbf{Pre-conditions} & The patient is signed in and can view their own adherence records. \\ \hline
\textbf{Post-conditions} & \textbf{Success:} Permitted adherence history is displayed without modifying it.\par
\textbf{Failure:} No record is changed. \\ \hline
\end{longtable}
% ---------- 05. Records and history ----------
% UC-10 - View Medication and Safety History | from UC-06
\clearpage
\subsection{Use-case: View Medication and Safety History}
\label{uc:uc-10}
\InlineUCModel{10}
\begin{longtable}{|p{0.20\textwidth}|p{0.71\textwidth}|}
\hline
\endfirsthead
\multicolumn{2}{l}{\textbf{UC-10: Specification (continued)}} \\
\hline
\endhead
\textbf{Use case Name} & UC-10: View Medication and Safety History \\ \hline
\textbf{Brief description} & The patient reviews saved medication records, past safety results, review outcomes, and their dates. \\ \hline
\textbf{Actors} & Patient / General User \\ \hline
\textbf{Basic Flow} &
1. The patient opens Medication and Safety History.\par
2. The system lists permitted medication records and past safety assessments.\par
3. The patient selects a record or assessment.\par
4. The system shows its date, findings, limitations, and linked HCP review status.\par
5. The system clearly labels results that no longer match the current medication record.
\\ \hline
\textbf{Alternative Flows} &
\textbf{A1. No history is available (step 2)}\par
1. The system displays an empty-state message and ends the use case.
\\ \hline
\textbf{Alternative Flows} &
\textbf{A2. A record is no longer available or access is denied (step 3)}\par
1. The system reports the limitation without showing private data.\par
2. The patient may select another permitted record.
\\ \hline
\textbf{Pre-conditions} & The patient is signed in and can view their own records. \\ \hline
\textbf{Post-conditions} & \textbf{Success:} Historical information is displayed without changing stored data.\par
\textbf{Failure:} No record or treatment is changed. \\ \hline
\end{longtable}
% ============================================================
% 06. CHATBOT AND PATIENT HEALTH SUMMARY
% ============================================================
% UC-11 - Ask a Medication Question
\clearpage
\subsection{Use-case: Ask a Medication Question}
\label{uc:uc-11}
\InlineUCModel{11}
\begin{longtable}{|p{0.20\textwidth}|p{0.71\textwidth}|}
\hline
\endfirsthead
\multicolumn{2}{l}{\textbf{UC-11: Specification (continued)}}\\
\hline
\endhead
\textbf{Use case Name} & UC-11: Ask a Medication Question \\ \hline
\textbf{Brief description} &
The patient receives a supported medication-related answer from the chatbot with appropriate limits.
\\ \hline
\textbf{Actors} & Patient / General User \\ \hline
\textbf{Basic Flow} &
\textcolor{blue}{1. The patient enters a medication-related question in the chatbot.}\par
\textcolor{blue}{2. The system identifies the question and retrieves relevant approved reference information and permitted patient context.}\par
\textcolor{blue}{3. The system prepares an answer with supporting information and relevant limitations.}\par
\textcolor{blue}{4. The system displays the answer or asks a follow-up question when clarification is needed.}\par
\textcolor{blue}{5. The patient may open View Conversation History to review previous conversations.}
\\ \hline
\textbf{Alternative Flows} &
\textcolor{blue}{\textbf{A1. Question is unclear (step 2)}}\par
\textcolor{blue}{1. The system asks a short follow-up question.}\par
\textcolor{blue}{2. If clarified, return to step 2; otherwise, end the use case.}
\\ \hline
\textbf{Alternative Flows} &
\textcolor{blue}{\textbf{A2. Patient requests source details (step 3)}}\par
\textcolor{blue}{1. The system displays the available supporting sources through View Supporting References.}\par
\textcolor{blue}{2. The patient may return to the answer.}
\\ \hline
\textbf{Alternative Flows} &
\textcolor{blue}{\textbf{A3. Reliable information is unavailable (step 3)}}\par
\textcolor{blue}{1. The system explains the information limit instead of generating unsupported medical claims.}
\\ \hline
\textbf{Alternative Flows} &
\textcolor{blue}{\textbf{A4. Request requires a clinical decision or another workflow (step 2)}}\par
\textcolor{blue}{1. The chatbot does not diagnose, prescribe, or change treatment.}\par
\textcolor{blue}{2. The system directs the patient to the appropriate supported workflow or qualified HCP when necessary.}\par
\textcolor{blue}{3. No medication record is changed solely from an unconfirmed chatbot message.}
\\ \hline
\textbf{Pre-conditions} & The patient is signed in and the chatbot is available. \\ \hline
\textbf{Post-conditions} & \textbf{Success:} The patient receives a supported answer, clarification, or clear limitation.\par
\textbf{Failure:} No unsupported clinical decision or unconfirmed record change occurs. \\ \hline
\end{longtable}
% UC-12 - Generate Patient Health Summary
\clearpage
\subsection{Use-case: Generate Patient Health Summary}
\label{uc:uc-12}
\InlineUCModel{12}
\begin{longtable}{|p{0.20\textwidth}|p{0.71\textwidth}|}
\hline
\endfirsthead
\multicolumn{2}{l}{\textbf{UC-12: Specification (continued)}}\\
\hline
\endhead
\textbf{Use case Name} & UC-12: Generate Patient Health Summary \\ \hline
\textbf{Brief description} &
\textcolor{blue}{The patient generates a structured summary from confirmed medication records and patient health information for review and optional sharing.}
\\ \hline
\textbf{Actors} & Patient / General User \\ \hline
\textbf{Basic Flow} &
\textcolor{blue}{1. The patient requests a Patient Health Summary.}\par
\textcolor{blue}{2. The system retrieves confirmed medication records, health conditions, allergies, and available safety assessment results.}\par
\textcolor{blue}{3. The system prepares a draft summary with information sources, dates, medication status, and known information gaps.}\par
\textcolor{blue}{4. The patient reviews the draft through Review and Confirm Draft Summary and corrects inaccurate information when necessary.}\par
\textcolor{blue}{5. The patient confirms the summary, and the system saves it for later viewing or optional sharing.}
\\ \hline
\textbf{Alternative Flows} &
\textcolor{blue}{\textbf{A1. Confirmed information is incomplete (step 2)}}\par
\textcolor{blue}{1. The system identifies missing information and prepares a limited draft using available confirmed data.}\par
\textcolor{blue}{2. If no eligible information is available, the system asks the patient to provide or confirm the required information.}
\\ \hline
\textbf{Alternative Flows} &
\textcolor{blue}{\textbf{A2. The draft needs correction (step 4)}}\par
\textcolor{blue}{1. The patient corrects the relevant information through the appropriate workflow.}\par
\textcolor{blue}{2. The system updates the draft and returns to step 4.}
\\ \hline
\textbf{Alternative Flows} &
\textbf{A3. Patient cancels or saving fails (step 5)}\par
1. No new confirmed summary is saved or shared.\par
2. The patient may retry if saving failed.
\\ \hline
\textbf{Pre-conditions} &
\textcolor{blue}{The patient is signed in and has at least one confirmed, permitted medication record or patient health detail.}
\\ \hline
\textbf{Post-conditions} &
\textbf{Success:} A patient-reviewed summary is stored for later viewing or optional sharing; it is not a diagnosis.\par
\textbf{Failure:} No unconfirmed summary is shared.
\\ \hline
\end{longtable}
% UC-13 - Share Patient Health Summary with HCP
\clearpage
\subsection{Use-case: Share Patient Health Summary with HCP}
\label{uc:uc-13}
\InlineUCModel{13}
\begin{longtable}{|p{0.20\textwidth}|p{0.71\textwidth}|}
\hline
\endfirsthead
\multicolumn{2}{l}{\textbf{UC-13: Specification (continued)}}\\
\hline
\endhead
\textbf{Use case Name} & UC-13: Share Patient Health Summary with HCP \\ \hline
\textbf{Brief description} &
The patient grants an authorized HCP access to a selected, confirmed health summary.
\\ \hline
\textbf{Actors} & Patient / General User \\ \hline
\textbf{Basic Flow} &
\textcolor{blue}{1. The patient selects a saved, confirmed health summary and chooses to share it with an HCP.}\par
\textcolor{blue}{2. The system displays the summary, its date, and the available authorized HCP recipients.}\par
\textcolor{blue}{3. The patient selects an HCP and chooses which information to share.}\par
\textcolor{blue}{4. The system displays the selected sharing scope and intended recipient.}\par
\textcolor{blue}{5. The patient reviews the sharing scope and confirms permission through Review Sharing Scope and Confirm Permission.}\par
\textcolor{blue}{6. The system validates the recipient and grants access only to the selected summary information.}\par
\textcolor{blue}{7. The system displays the sharing status. The patient may later use Revoke Summary Sharing to withdraw access.}
\\ \hline
\textbf{Alternative Flows} &
\textbf{A1. Patient declines permission (step 5)}\par
1. No new access is granted; the use case ends.
\\ \hline
\textbf{Alternative Flows} &
\textcolor{blue}{\textbf{A2. Recipient is unauthorized (step 6)}}\par
\textcolor{blue}{1. The system denies sharing and informs the patient.}\par
\textcolor{blue}{2. The patient may select another authorized HCP or cancel.}
\\ \hline
\textbf{Alternative Flows} &
\textcolor{blue}{\textbf{A3. Summary is unconfirmed or unavailable (steps 1--2)}}\par
\textcolor{blue}{1. The system does not allow the summary to be shared.}\par
\textcolor{blue}{2. The patient may return to Generate Patient Health Summary (UC-12).}
\\ \hline
\textbf{Alternative Flows} &
\textbf{A4. Sharing fails (step 6)}\par
1. The system reports the error and does not claim success or create unintended access.
\\ \hline
\textbf{Pre-conditions} &
The patient is signed in and has a saved, reviewed summary; an authorized target can be selected.
\\ \hline
\textbf{Post-conditions} &
\textcolor{blue}{\textbf{Success:} Only the selected summary information is accessible to the authorized HCP. The patient can revoke this permission later. HCP case-sharing permission remains separate from summary-sharing permission.}\par
\textbf{Failure:} No new access is granted.
\\ \hline
\end{longtable}
% ============================================================
% 07. DATASET REVIEW
% ============================================================
% UC-14 - Review Medication Reference Data
\clearpage
\subsection{Use-case: Review Medication Reference Data}
\label{uc:uc-14}
\InlineUCModel{14}
\begin{longtable}{|p{0.20\textwidth}|p{0.71\textwidth}|}
\hline
\endfirsthead
\multicolumn{2}{l}{\textbf{UC-14: Specification (continued)}}\\
\hline
\endhead
\textbf{Use case Name} &
\textcolor{blue}{UC-14: Review Medication Reference Data}
\\ \hline
\textbf{Brief description} &
\textcolor{blue}{The Dataset Reviewer checks medication reference data against source evidence and records a traceable data-quality decision.}
\\ \hline
\textbf{Actors} & Dataset Reviewer \\ \hline
\textbf{Basic Flow} &
\textcolor{blue}{1. The reviewer opens the permitted medication reference data list and selects a record.}\par
\textcolor{blue}{2. The system displays the record, ingredients, mappings, source references, dataset version, and review status.}\par
\textcolor{blue}{3. The reviewer uses Compare Record with Evidence to compare the record with available source information.}\par
\textcolor{blue}{4. The reviewer uses Record Data-Quality Review Decision to select Accept, Reject, or Needs Correction and provide a reason.}\par
\textcolor{blue}{5. The system validates the review details and saves the decision, reviewer, time, and supporting evidence.}\par
6. The system updates the dataset review status.\par
\textcolor{blue}{7. The reviewer may open View Dataset Review History to inspect previous decisions.}
\\ \hline
\textbf{Alternative Flows} &
\textcolor{blue}{\textbf{A1. Source evidence is missing or inconsistent (steps 2--4)}}\par
\textcolor{blue}{1. The reviewer selects Needs Correction or Reject and records the issue.}\par
\textcolor{blue}{2. The system does not mark the record as clinically validated.}\par
\textcolor{blue}{3. Continue to step 5 to save the review decision.}
\\ \hline
\textbf{Alternative Flows} &
\textbf{A2. Record changed during review (step 5)}\par
1. The system asks the reviewer to reload the latest version before deciding.\par
2. Return to step 2.
\\ \hline
\textbf{Alternative Flows} &
\textbf{A3. Reviewer cancels or saving fails (steps 4--5)}\par
1. No new decision is saved; the previous review status remains.
\\ \hline
\textbf{Pre-conditions} &
The Dataset Reviewer is signed in and permitted to review the record.
\\ \hline
\textbf{Post-conditions} &
\textbf{Success:} A source-linked data-quality decision and change history are recorded; no independent clinical validation is claimed.\par
\textbf{Failure:} No new review decision is saved.
\\ \hline
\end{longtable}
% ============================================================
% 08. ADMINISTRATION
% ============================================================
% UC-15 - Manage User Accounts and Roles
\clearpage
\subsection{Use-case: Manage User Accounts and Roles}
\label{uc:uc-15}
\InlineUCModel{15}
\begin{longtable}{|p{0.20\textwidth}|p{0.71\textwidth}|}
\hline
\endfirsthead
\multicolumn{2}{l}{\textbf{UC-15: Specification (continued)}}\\
\hline
\endhead
\textbf{Use case Name} & UC-15: Manage User Accounts and Roles \\ \hline
\textbf{Brief description} &
The Administrator creates or updates accounts and assigns permitted roles under the system access policy.
\\ \hline
\textbf{Actors} & Administrator \\ \hline
\textbf{Basic Flow} &
\textcolor{blue}{1. The Administrator opens Manage User Accounts and Roles.}\par
\textcolor{blue}{2. The system displays the account list, status, and permitted role options.}\par
\textcolor{blue}{3. The Administrator selects an account-management operation.}\par
4. The system validates the request against account and RBAC rules.\par
5. The Administrator confirms the proposed change.\par
6. The system saves the authorized change and records the administrative activity.\par
\textcolor{blue}{7. The Administrator may open View Account Activity History to inspect account-related activities.}
\\ \hline
\textbf{Alternative Flows} &
\textcolor{blue}{\textbf{A1. Administrator creates a user account (step 3)}}\par
\textcolor{blue}{1. The Administrator selects Create User Account and enters the required account details.}\par
\textcolor{blue}{2. Continue to step 4.}
\\ \hline
\textbf{Alternative Flows} &
\textcolor{blue}{\textbf{A2. Administrator assigns a user role (step 3)}}\par
\textcolor{blue}{1. The Administrator selects Assign User Role, chooses an account, and selects a permitted role.}\par
\textcolor{blue}{2. Continue to step 4.}
\\ \hline
\textbf{Alternative Flows} &
\textcolor{blue}{\textbf{A3. Administrator deactivates a user account (step 3)}}\par
\textcolor{blue}{1. The Administrator selects Deactivate User Account and chooses the target account.}\par
\textcolor{blue}{2. The system deactivates account access without deleting the account's records.}\par
\textcolor{blue}{3. Continue to step 5 for confirmation.}
\\ \hline
\textbf{Alternative Flows} &
\textcolor{blue}{\textbf{A4. Administrator resets a user password (step 3)}}\par
\textcolor{blue}{1. The Administrator selects Reset User Password and chooses the target account.}\par
\textcolor{blue}{2. The system prepares the approved password-reset process.}\par
\textcolor{blue}{3. Continue to step 4.}
\\ \hline
\textbf{Alternative Flows} &
\textcolor{blue}{\textbf{A5. Request is invalid or not permitted (step 4)}}\par
\textcolor{blue}{1. The system rejects the request without applying the change or granting unauthorized permissions.}\par
\textcolor{blue}{2. The Administrator may correct the request or select another permitted operation.}
\\ \hline
\textbf{Alternative Flows} &
\textbf{A6. Administrator cancels or save fails (steps 5--6)}\par
1. The previous account and role settings remain unchanged.\par
2. The system reports a save error when relevant.
\\ \hline
\textbf{Pre-conditions} &
The Administrator is signed in with account and role-management permission.
\\ \hline
\textbf{Post-conditions} &
\textbf{Success:} The permitted account or role change is stored with an activity record.\par
\textbf{Failure:} No unauthorized change is saved.
\\ \hline
\end{longtable}
% ============================================================
% 09. TECHNICAL SUPPORT
% ============================================================
% UC-16 - Report Technical Issue
\clearpage
\subsection{Use-case: Report Technical Issue}
\label{uc:uc-16}
\InlineUCModel{16}
\begin{longtable}{|p{0.20\textwidth}|p{0.71\textwidth}|}
\hline
\endfirsthead
\multicolumn{2}{l}{\textbf{UC-16: Specification (continued)}}\\
\hline
\endhead
\textbf{Use case Name} &
\textcolor{blue}{UC-16: Report Technical Issue}
\\ \hline
\textbf{Brief description} &
A registered user reports a technical problem and receives a support ticket status.
\\ \hline
\textbf{Actors} &
\textcolor{blue}{Registered User (all registered roles)}
\\ \hline
\textbf{Basic Flow} &
1. The user opens Report Technical Issue.\par
2. The system shows the required issue fields and optional diagnostic details.\par
3. The user describes the problem and reviews any attached information.\par
\textcolor{blue}{4. The system displays the report details through Review Issue Details and Attachments.}\par
5. The user submits the report.\par
\textcolor{blue}{6. The system validates the required information and attachment rules.}\par
\textcolor{blue}{7. The system saves the ticket and displays its identifier and initial status only after successful creation.}
\\ \hline
\textbf{Alternative Flows} &
\textbf{A1. Required information is missing (step 6)}\par
1. The system highlights the missing information.\par
2. Return to step 3 for correction.
\\ \hline
\textbf{Alternative Flows} &
\textbf{A2. Attachment contains unnecessary patient information (step 3)}\par
1. The system requests removal or masking of the private data before submission.\par
2. Return to step 3.
\\ \hline
\textbf{Alternative Flows} &
\textbf{A3. User cancels or ticket creation fails (steps 5--7)}\par
1. No successful ticket submission is reported.\par
2. The user may retry after an error or end the use case.
\\ \hline
\textbf{Pre-conditions} &
The user is signed in and can open the technical support feature.
\\ \hline
\textbf{Post-conditions} &
\textbf{Success:} A support ticket is stored and its identifier and initial status are displayed without unnecessary patient data.\par
\textbf{Failure:} No ticket is created.
\\ \hline
\end{longtable}
% UC-17 - Handle Technical Support Request
\clearpage
\subsection{Use-case: Handle Technical Support Request}
\label{uc:uc-17}
\InlineUCModel{17}
\begin{longtable}{|p{0.20\textwidth}|p{0.71\textwidth}|}
\hline
\endfirsthead
\multicolumn{2}{l}{\textbf{UC-17: Specification (continued)}}\\
\hline
\endhead
\textbf{Use case Name} &
\textcolor{blue}{UC-17: Handle Technical Support Request}
\\ \hline
\textbf{Brief description} &
Technical Support investigates a permitted ticket, records progress, and provides an outcome when available.
\\ \hline
\textbf{Actors} & Technical Support \\ \hline
\textbf{Basic Flow} &
\textcolor{blue}{1. Support staff opens View Assigned Support Tickets and selects a ticket.}\par
2. The system shows the issue description, status, and allowed technical details.\par
3. Support staff investigates the issue and records relevant notes.\par
4. Support staff records a resolution or updates the ticket status.\par
\textcolor{blue}{5. The system saves the action, time, and responsible account and displays the updated ticket status to the reporting user.}
\\ \hline
\textbf{Alternative Flows} &
\textbf{A1. More information is required (step 3)}\par
1. Support staff requests only necessary technical details through Request Additional Technical Information.\par
2. The system records the request and marks the ticket Pending User Response.\par
3. The reporting user may provide additional information through the same ticket.\par
4. The system records the reply and returns the ticket to the support queue for continued handling.
\\ \hline
\textbf{Alternative Flows} &
\textbf{A2. The issue is not resolved (step 4)}\par
1. Support staff records progress and keeps the ticket Open or In Progress.\par
2. Continue to step 5.
\\ \hline
\textbf{Alternative Flows} &
\textbf{A3. Unauthorized patient data is present (step 2)}\par
1. The system restricts access to medical content.\par
2. Support staff works with permitted technical details only.
\\ \hline
\textbf{Alternative Flows} &
\textbf{A4. Access is denied or saving fails (steps 1 or 5)}\par
1. The system denies the action or reports the error.\par
2. No unsupported resolution is recorded.
\\ \hline
\textbf{Pre-conditions} &
Support staff is signed in and permitted to access the selected support ticket.
\\ \hline
\textbf{Post-conditions} &
\textbf{Success:} The ticket status and troubleshooting history are saved; the reporting user can view the updated status.\par
\textbf{Failure:} No unsupported ticket update or patient-data access occurs.
\\ \hline
\end{longtable}
% ============================================================
% 10. PATIENT HEALTH PROFILE
% ============================================================
% UC-18 - Manage Patient Health Profile
\clearpage
\subsection{Use-case: Manage Patient Health Profile}
\label{uc:uc-18}
\InlineUCModel{18}
\begin{longtable}{|p{0.20\textwidth}|p{0.71\textwidth}|}
\hline
\endfirsthead
\multicolumn{2}{l}{\textbf{UC-18: Specification (continued)}}\\
\hline
\endhead
\textbf{Use case Name} &
UC-18: Manage Patient Health Profile
\\ \hline
\textbf{Brief description} &
\textcolor{blue}{The patient manages personal health information, including medical conditions, allergies, and medication history, for use in supported system workflows.}
\\ \hline
\textbf{Actors} & Patient / General User \\ \hline
\textbf{Basic Flow} &
\textcolor{blue}{1. The patient opens Manage Patient Health Profile.}\par
\textcolor{blue}{2. The system displays existing medical conditions, allergies, and medication history with their available sources and update information.}\par
\textcolor{blue}{3. The patient selects a category and adds or updates the relevant health information.}\par
\textcolor{blue}{4. The system validates the required information and prepares the changes for review.}\par
\textcolor{blue}{5. The patient reviews and confirms the changes through Review and Confirm Health Information.}\par
\textcolor{blue}{6. The system saves the confirmed information together with its source, status, and update time.}
\\ \hline
\textbf{Alternative Flows} &
\textcolor{blue}{\textbf{A1. Required information is incomplete or invalid (step 4)}}\par
\textcolor{blue}{1. The system identifies the missing or invalid information.}\par
\textcolor{blue}{2. The patient corrects the information and returns to step 4.}
\\ \hline
\textbf{Alternative Flows} &
\textcolor{blue}{\textbf{A2. Patient cancels the changes (step 5)}}\par
\textcolor{blue}{1. The system discards the unconfirmed changes.}\par
\textcolor{blue}{2. The previously saved information remains unchanged.}
\\ \hline
\textbf{Pre-conditions} &
\textcolor{blue}{The patient is signed in and can access their own health profile.}
\\ \hline
\textbf{Post-conditions} &
\textcolor{blue}{\textbf{Success:} Confirmed health information is stored with its available source, status, and update time. Patient-provided information is not automatically treated as clinically validated.}\par
\textcolor{blue}{\textbf{Failure:} No unconfirmed changes are saved as confirmed health information.}
\\ \hline
\end{longtable}
% ============================================================
% 11. PERSONAL MEDICATION CABINET
% ============================================================
% UC-19 - Manage Personal Medication Cabinet
\clearpage
\subsection{Use-case: Manage Personal Medication Cabinet}
\label{uc:uc-19}
\InlineUCModel{19}
\begin{longtable}{|p{0.20\textwidth}|p{0.71\textwidth}|}
\hline
\textbf{Use case Name} &
\textcolor{blue}{UC-19: Manage Personal Medication Cabinet}
\\ \hline
\textbf{Brief description} &
\textcolor{blue}{The patient manages a personal collection of medication records, including medication details and usage status. This use case is reserved for future development.}
\\ \hline
\textbf{Actors} &
Patient / General User
\\ \hline
\textbf{Basic Flow} &
\textcolor{blue}{1. The patient opens Manage Personal Medication Cabinet.}\par
\textcolor{blue}{2. The system displays the medications saved in the cabinet.}\par
\textcolor{blue}{3. The patient selects an operation: add a medication, view its details, update its information, update its usage status, or view expired and expiring medications.}\par
\textcolor{blue}{4. The system displays the relevant information or input fields.}\par
\textcolor{blue}{5. The patient confirms the requested changes.}\par
\textcolor{blue}{6. The system validates and saves the changes, then displays the updated information.}
\\ \hline
\textbf{Alternative Flows} &
\textcolor{blue}{A1. Required information is missing or invalid (step 5)}\par
\textcolor{blue}{1. The system highlights the missing or invalid information.}\par
\textcolor{blue}{2. The patient corrects the information and retries.}
\\ \hline
\textbf{Pre-conditions} &
\textcolor{blue}{The patient is signed in and can access the personal medication cabinet.}
\\ \hline
\textbf{Post-conditions} &
\textcolor{blue}{Success: Medication information and usage status are updated in the cabinet.}\par
\textcolor{blue}{Failure: No invalid or unconfirmed changes are saved.}
\\ \hline
\end{longtable}
\textcolor{blue}{\textbf{Scope note: This use case is reserved for future development and is not included in the current implementation scope.}}
% ============================================================
% 12. PERSONAL ACCOUNT
% ============================================================
% UC-20 - Manage Personal Account
\clearpage
\subsection{Use-case: Manage Personal Account}
\label{uc:uc-20}
\InlineUCModel{20}
\begin{longtable}{|p{0.20\textwidth}|p{0.71\textwidth}|}
\hline
\textbf{Use case Name} &
UC-20: Manage Personal Account
\\ \hline
\textbf{Brief description} &
A registered user views and updates personal account information or changes their password.
\\ \hline
\textbf{Actors} &
Registered User (Patient / HCP)
\\ \hline
\textbf{Basic Flow} &
\textcolor{blue}{1. The user opens Manage Personal Account.}\par
\textcolor{blue}{2. The system displays the current account information and available operations.}\par
\textcolor{blue}{3. The user selects View Personal Account, Update Personal Account Details, or Change Password.}\par
\textcolor{blue}{4. The system performs the selected operation according to the user's permissions.}\par
\textcolor{blue}{5. The system displays the result and saves any successful changes.}
\\ \hline
\textbf{Alternative Flows} &
\textbf{A1. View personal account (step 3)}\par
1. The system displays the user's account information, such as name and email.\par
2. No account information is changed.
\\ \hline
\textbf{Alternative Flows} &
\textbf{A2. Update personal account details (step 3)}\par
1. The user edits the permitted personal information.\par
2. The system validates and saves the changes.\par
3. The system displays the updated information.
\\ \hline
\textbf{Alternative Flows} &
\textbf{A3. Change password (step 3)}\par
1. The user enters the current password and a new password with confirmation.\par
2. The system verifies the current password and validates the new password.\par
3. The system updates the password and confirms the change.
\par
\textcolor{blue}{\textbf{A4. Account information or password is invalid}}\par
\textcolor{blue}{1. The system displays an error message.}\par
\textcolor{blue}{2. No invalid changes are saved.}\par
\textcolor{blue}{3. The user may correct the information and retry.}
\\ \hline
\textbf{Pre-conditions} &
The user is signed in and can access their own personal account.
\\ \hline
\textbf{Post-conditions} &
\textbf{Success:} The requested account information or password is updated. The user's assigned role and access permissions remain unchanged.\par
\textbf{Failure:} No unauthorized or invalid account changes are saved.
\\ \hline
\end{longtable}